\documentclass[12pt]{article}
\usepackage[utf8]{inputenc}
\usepackage{amsmath, amssymb, geometry, graphicx, setspace, xcolor, listings}
\geometry{letterpaper, margin=1in}
\setstretch{1.15}

% Code Block Styling
\lstset{
    basicstyle=\ttfamily\small,
    keywordstyle=\color{blue}\bfseries,
    commentstyle=\color{green!50!black},
    frame=single,
    backgroundcolor=\color{gray!10}
}

\title{\textbf{GMAT Validation Guide: \\ From Lambert's Problem to 3D Departure}}
\author{Module: Astrodynamics 2026-1}
\date{}

\begin{document}

\maketitle

\section*{Objective}
The Python SPICE script used Lambert's problem to find the optimal launch date and the required Heliocentric Velocity at departure ($V_{tx1}$). In this lab, you will extract the Hyperbolic Excess Velocity vector ($\mathbf{V}_\infty$) from those results, calculate the analytical geometry ($\beta$ and $\Delta V$), and use GMAT's Differential Corrector to simulate the exact 3D departure hyperbola.

\section*{Step 1: Translating Python Vectors to GMAT Targets}
GMAT cannot target a heliocentric vector directly while the spacecraft is still in Earth orbit. Instead, it targets the energy and 3D direction of the escape asymptote. 

Using the output from your Python script, calculate the $\mathbf{V}_\infty$ vector:
\begin{equation}
    \mathbf{V}_\infty = \mathbf{V}_{tx1} - \mathbf{V}_{Earth} = \begin{bmatrix} V_X \\ V_Y \\ V_Z \end{bmatrix}
\end{equation}

You must convert this Cartesian vector into the three parameters GMAT uses for interplanetary targeting:
\begin{enumerate}
    \item \textbf{Characteristic Energy ($C_3$):} The squared magnitude of the vector.
    \begin{equation}
        C_3 = |\mathbf{V}_\infty|^2 = V_X^2 + V_Y^2 + V_Z^2
    \end{equation}
    \item \textbf{Right Ascension of the Asymptote (RLA):} The angle in the XY plane.
    \begin{equation}
        RLA = \arctan2(V_Y, V_X) \times \left(\frac{180}{\pi}\right)
    \end{equation}
    \item \textbf{Declination of the Asymptote (DLA):} The angle above/below the equator.
    \begin{equation}
        DLA = \arcsin\left(\frac{V_Z}{|\mathbf{V}_\infty|}\right) \times \left(\frac{180}{\pi}\right)
    \end{equation}
\end{enumerate}

\section*{Step 2: Calculating the Analytical Initial Guesses}
The Differential Corrector will fail if your starting guesses are completely wrong. Assume a circular parking orbit altitude of $h = 300$ km ($r_p = 6678.136$ km). 

Calculate your initial guesses using the 2-body geometry learned in class:
\begin{itemize}
    \item \textbf{Guess for the Burn ($\Delta V$):}
    \begin{equation*}
        \Delta V = \sqrt{C_3 + \frac{2\mu_\oplus}{r_p}} - \sqrt{\frac{\mu_\oplus}{r_p}}
    \end{equation*}
    \item \textbf{Guess for the Burn Location ($\beta$):} \\
    Calculate the eccentricity ($e = 1 + r_p C_3 / \mu_\oplus$). The Asymptote Angle ($\beta = \cos^{-1}(1/e)$) tells you approximately how far away from the target RLA the True Anomaly (TA) should be.
\end{itemize}

\newpage

\section*{Step 3: NASA GMAT Setup}
Create a new mission in GMAT. Ensure the gravity model is set to \textbf{Point Mass} and all drag/solar radiation pressure models are disabled to perfectly match your analytical Lambert assumptions.

\subsection*{1. Spacecraft and Maneuver Setup}
\begin{itemize}
    \item \textbf{DefaultSC:} Set to Keplerian elements.
    \item \textbf{SMA:} 6678.136 km
    \item \textbf{ECC:} 0.0
    \item \textbf{INC:} 28.5$^\circ$ (Standard Kennedy Space Center launch inclination)
    \item \textbf{TA:} Set to your $\beta$ angle guess.
    \item Create an \textbf{ImpulsiveBurn} in the VNB (Velocity-Normal-Binormal) frame.
\end{itemize}

\subsection*{2. The Mission Sequence (Targeting)}
Because the target vector ($\mathbf{V}_\infty$) exists in 3D space, you must provide GMAT with three control variables to hit three targets.

Create a \texttt{Target} sequence and configure the Differential Corrector as follows:

\begin{lstlisting}
% --- 1. The Control Variables (Vary) ---
% Vary the Prograde Burn Magnitude
Vary DC1(ImpulsiveBurn1.Element1 = [Your Analytical \Delta V Guess]);

% Vary the Burn Location in the Orbit
Vary DC1(DefaultSC.TA = [Your Analytical \beta Guess]);

% Vary the Orientation of the Orbital Plane
Vary DC1(DefaultSC.RAAN = 0.0, {Perturbation = 0.0001, MaxStep = 5});

% --- 2. Apply and Propagate ---
Maneuver ImpulsiveBurn1(DefaultSC);
Propagate DefaultProp(DefaultSC) {DefaultSC.Earth.RMAG = 925000};

% --- 3. The Goals (Achieve) ---
% Target the Energy
Achieve DC1(DefaultSC.Earth.C3Energy = [Your Calculated C3]);

% Target the 3D Direction
Achieve DC1(DefaultSC.Earth.RLA = [Your Calculated RLA]);
Achieve DC1(DefaultSC.Earth.DLA = [Your Calculated DLA]);
\end{lstlisting}

\section*{Step 4: Lab Deliverables}
Run the sequence. The Differential Corrector will iteratively adjust the burn magnitude ($\Delta V$), the burn location (TA), and the orbital plane (RAAN) until the spacecraft's escape asymptote perfectly matches the Lambert solution from your Python script.

\textbf{Submit the following:}
\begin{enumerate}
    \item Your Python output displaying the optimal Heliocentric vector.
    \item Your hand calculations for $C_3$, RLA, DLA, $\Delta V$, and $\beta$.
    \item A screenshot of the converged GMAT Message Window showing the final $\Delta V$ and TA values. Briefly explain why GMAT's final True Anomaly differs slightly from your exact $\beta$ calculation (Hint: Think about 3D spherical geometry vs. 2D planar geometry).
    \item A screenshot of the 3D escape hyperbola from GMAT's OrbitView.
\end{enumerate}

\end{document}